% !PRINTED 244
\weberpage{244}

Auch der Fall $D=0$ giebt hier zu keinen weiteren Unterscheidungen Anlass; denn in diesem Falle sind die beiden Wurzeln der quadratischen Gleichung einander gleich und reell, von den drei Wurzeln der cubischen Gleichung zwei einander gleich und alle drei reell; denn eine imaginäre Wurzel kann hier nicht doppelt vorkommen, weil sonst auch die conjugirte doppelt vorkommen würde, und ebenso wenig kann die einzelne Wurzel imaginär sein, weil sonst eine zweite vorhanden sein müsste.

Es kann sich bei der cubischen Gleichung nur noch darum handeln, die Bedingung dafür aufzusuchen, dass alle drei Wurzeln einander gleich sind.

In diesem Falle muss die linke Seite der cubischen Gleichung ein vollständiger Cubus sein, also
\[
a_0 x^3 + a_1 x^2 + a_2 x + a_3 = a_0 (x - \alpha)^3.
\]
Die Vergleichung beider Seiten dieser identischen Gleichung ergiebt
\[
a_1 = -3a_0\alpha, \quad a_2 = 3a_0\alpha^2, \quad a_3 = -a_0\alpha^3,
\]
woraus man durch Elimination von $\alpha$ erhält
\weq{1}{\begin{aligned} a_1^2 - 3a_0a_2 &= 0,\\ a_1a_2 - 9a_0a_3 &= 0. \end{aligned}}
Daraus ergiebt sich als Folge (indem man die erste dieser Gleichungen mit $a_2$, die zweite mit $a_1$ multiplicirt und subtrahirt)
\[
a_2^2 - 3a_1a_3 = 0;
\]
und wenn diese Bedingungen erfüllt sind, so folgt daraus umgekehrt
\[
f(x) = a_0\left(x + \frac{a_1}{3a_0}\right)^3,
\]
d.\,h. die Gleichheit aller drei Wurzeln.

Bei der cubischen Gleichung
\[
x^3 + a_1 x^2 + a_2 x + a_3 = 0
\]
können wir ausser über die Realität auch noch über die Vorzeichen der Wurzeln vollständig entscheiden.

Bezeichnen wir nämlich mit $\alpha$, $\beta$, $\gamma$ die drei Wurzeln, so ist
\weq{2}{a_1 = -(\alpha+\beta+\gamma), \quad a_2 = \alpha\beta+\alpha\gamma+\beta\gamma, \quad a_3 = -\alpha\beta\gamma.}

Ist nun $D<0$, also nur eine Wurzel, etwa $\alpha$, reell, so ist $\beta\gamma$ positiv und $\alpha$ wird negativ oder positiv sein, je nachdem $a_3$ positiv oder negativ ist. Ist aber $D$ positiv, also alle drei Wurzeln reell, so ist, wenn $\alpha$, $\beta$, $\gamma$ positiv sind, nach (2)
%CONT>>
